\scanpage{91}
% archive sheet number: 93

% slanted
\margin{\uline{Déc 83}}

\uline{Courbes elliptiques réelles}

Soit $k$ un anneau (\unsure{présumé|disons|par exemple} comm.) \deleted{$M$ un $k$-\unsure{mo}} Considérons la catégorie des $k$-Modules $M$ munis d'un autom $\sigma$ \deleted{tel q} tels que

a) $\sigma^2=\mathrm{id}$

b) $2\mathrm{id}_M$ injectif (on dit alors $M$ 2-régulier)

A un tel couple $(M,\sigma)$, \inserted{nous} associons \deleted{\unsure{les}} couples $(P,Q,m)$ de deux \inserted{$k$-}modules pour lesquels $2\mathrm{id}_P$ et $2\mathrm{id}_Q$ sont injectifs (i.e. $P,Q$ 2-réguliers), et $m\subset P/2P\oplus Q/2Q$ \inserted{sous $k/2k$-module} avec $m\cap P/2P=0$, $m\cap Q/2Q=0$ (soit deux sous-modules $p'\subset p\overset{\mathrm{déf}}=P/2P$, $q'\subset q=Q/2Q$, et un iso. $p'\simeq q'$), de la façon suivante

\[
\begin{aligned}
P&=M_+=\operatorname{Ker}(\mathrm{id}_M-\sigma)=\{x\in M\mid\sigma x=x\}\\
Q&=M_-=\operatorname{Ker}(\mathrm{id}_M-\sigma)=\{\unsure{n|x}\in M\mid\sigma x=-x\}
\end{aligned}
\]

\unsure{NB|N.B.} $M_+\cap M_-=0$ dans $M$, car $M$ 2-régulier, et on \margin{!} a $2M\subset M_+\oplus M_-\hookrightarrow M\deleted{\subset\tfrac12(M_+\oplus M_-)}$ car $\forall x\in M$, on a $2x=(\underset{M_+}{\underbrace{x+\sigma x}})+(\underset{M_-}{\underbrace{x-\sigma x}})$

soit
\[
\underset{\substack{\shortparallel\\2M}}{M'}\subset M_+\oplus M_{\unsure{1|-}}=P\oplus Q\qquad d'
\]

on a donc
\[
\deleted{\unsure{M'\hookrightarrow}}2P\oplus 2Q\subset M'\subset P\oplus Q
\]

donc $M'$ est l'image d'un sous-module $m$ de $(P\oplus Q)/2P\oplus 2Q\simeq\underset p{\underbrace{P/2P}}\oplus\underset q{\underbrace{Q/2Q}}$, et on vérifie \deleted{\unsure{que}}

\[
\begin{aligned}
\left\{\begin{aligned}m\cap\deleted{\unsure{q}}p&=0\\m\cap q&=0\end{aligned}\right.
&\quad\begin{aligned}
\text{i.e. } \underset{M'\cap P+(2P+2Q)}{\underbrace{M'\cap(P+(2P+2Q))}}&=(2P+2Q)\\
\text{i.e. }M'\cap(Q+(2P+2Q))&=2Q
\end{aligned}
\end{aligned}
\]
car $M'\cap P=2P$ puisque $2x\in P$ i.e. $2\sigma x=2x$ implique $x\in P$ i.e. $\sigma x=x$, et $M'\cap Q=2Q$.

D'ailleurs \deleted{\unsure{$M'\subset P\oplus Q$}} $(M,\sigma)$ se reconstitue à isom can près via $(P,Q,m)$, car $M'\subset P\oplus Q$ se reconstitue, et \deleted{\unsure{$M'$}}$\sigma_{M'}$ aussi bien entendu par $\mathrm{id}_P\oplus(-\mathrm{id}_Q)$), or $2\mathrm{id}_M\deleted{\unsure{4p}}:M\to M'$ est un iso de $(M,\sigma_M)$ avec $(M',\sigma_{M'})$.

\uline{Prop} Le foncteur précédent $(M,\sigma)\longmapsto(P,Q,m)$ est une équivalence

\uline{Cor} Supposons $k$ \deleted{\unsure{qui|soit}} \inserted{comm.} \deleted{intègre et \unsure{régulier} et ppal}\inserted{$I$ de Dedekind, avec \deleted{$k$}$k$ 2-régulier}, prenons \deleted{des} $M$ \deleted{libres} \inserted{projectifs} de \deleted{rang}\inserted{rg} 2 avec $\sigma\ne\mathrm{id}_M,-\mathrm{id}_M$, alors la cat \deleted{de ces}\inserted{des} foncteurs \deleted{équivaut à celle}\inserted{$(M,\sigma)$ avec, où $M$ fid\textsuperscript{t} de rg 2,} équivaut à celle des paires $(P,Q)$ avec $P,Q$ inv., et $p'\subset P/2P$, $q'\subset Q/2Q$ et $p'\overset u\simeq q'$ donné.
